\chapter{Lowering, optimisation and the emitted program's well-formedness}
\label{chap:lowering}

Chapters~\ref{chap:erases} and~\ref{chap:correctness} relate a \code{Lean.Expr} to a \lbox{}
term by $\Erases$ and prove that relation a forward simulation. That is not yet a statement
about what the shipping eraser writes to disk, because two of the eraser's steps are
\emph{compilations} that read the global environment rather than the term. Lean has no
primitive \code{case}: pattern matching goes through \code{casesOn} and \code{rec}
\emph{constants}, and the eraser recognises a saturated application of such a constant, drops
the parameters and the motive, peels each minor premise's lambda-chain into an alternative and
emits a \code{case} node. Lean's recursion is by \emph{environment reference}: a mutual block is
a set of constants naming one another, while \lbox{} recursion is a \code{fix} node with de
Bruijn-bound siblings, so the eraser erases each member with the block's names mapped to fresh
free variables and closes the result. Neither step is a relation between a \code{Lean.Expr} and
an \code{LBTerm} indexed by the source context alone. The development therefore factors the
eraser as a composite: an erasure into a \emph{specification} environment, in which eliminators
are still constants and block members still hold their plain bodies, followed by a
\lbox{}-to-\lbox{} \emph{pass}. This chapter is that pass --- the relation $\Lower$, its
metatheory, its inversion kit and the fixpoint round trip it needs --- together with everything
the development says about the program the pass produces. The Lean files are
\code{LeanToLambdaBox/Lower.lean}, \code{LowerFix.lean}, \code{FixUnfold.lean} and
\code{ElimBody.lean}.

The second half of the chapter is the output boundary, in \code{Output.lean},
\code{OutputShape.lean} and \code{OutputCheck.lean}. \code{LBWfPeregrine} is a twelve-clause
structure stating what peregrine's untyped transform pipeline needs of the emitted program: it
is \code{peregrine validate}'s checks, plus the constructor-saturation invariant that
\code{remove\_params\_optimization} consumes and \code{validate} omits, plus a printability
condition on binder names, each clause quantified over the emitted term \emph{and} every
constant body of the emitted environment through the combinator \code{OnProgram} --- the same
environment/term split as MetaRocq's \code{expanded\_eprogram\_cstrs}. What was deliberately
left out is as informative as what is in: MetaRocq's fixpoint eta-expansion invariant is
\emph{false} on this eraser's output, so it is stated separately as \code{LBExpandedFix}, the
true consumer precondition is named \code{PeregrinePre}, and the capstone concludes only the
weaker conjunct. \code{OutputCheck.lean} makes the twelve clauses decidable by a structurally
recursive Boolean checker, so that a concrete emitted program discharges them by kernel
computation; this is how each green rung supplies the capstone's well-formedness premise rather
than assuming it. The same discipline produces \code{NoBodylessRefs}, this repository's
decidable analogue of [S, \S7.3]'s \code{axiom\_free}, whose job is to keep a rung from being
vacuously green.

Standing slightly apart is \code{Optimize.lean}: a complete, \code{sorry}-free re-proof of
MetaRocq's \code{optimize} / \code{remove\_match\_on\_box} [S, \S7.4] as a genuine
\lbox{}-to-\lbox{} pass with a flag-discharging correctness theorem. It is the tree's declared
\emph{template} --- $\Lower$ is stated in its shape --- but it is not composed into the
capstone, it is stated at block-form flags rather than at the shipping applied-form flags, and
it lies outside the import closure of the capstone and of the rungs. The chapter says so
plainly rather than letting the surrounding documentation's ``worked example'' language suggest
that the pass runs.

\section{The pass relation \texorpdfstring{$\Lower$}{Lower}}
\label{sec:lowering:pass}

\begin{definition}[Environment queries of the pass]
\label{def:ElimDecl}
\lean{LeanToLambdaBox.DefnDecl, LeanToLambdaBox.IndBodyOf, LeanToLambdaBox.ElimDecl,
LeanToLambdaBox.RuntimeKey, LeanToLambdaBox.ClosedBodies, LeanToLambdaBox.FVarFreeBodies}
\leanok
\uses{def:ElimBody, def:GlobalDeclarations, def:LBTerm, def:Kername, def:InductiveId,
def:OneInductiveBody, def:LBTerm-envLookup, def:LBClosed, def:hasFVar}
Six predicates on a \emph{specification} environment $\Gamma$ --- the only thing about the
environment the pass ever reads. $\code{DefnDecl}\ \Gamma\ kn\ b$ says $kn$ is declared in
$\Gamma$ with body $b$. $\code{IndBodyOf}\ iid\ np\ nfs\ mib$ says the mutual block $mib$ has
$np$ parameters, that $iid$ selects a body of it, that the body is marked non-propositional,
and that its constructors' field counts are $nfs$ --- exactly the numbers the target semantics
reads when an iota step fires. $\code{ElimDecl}\ \Gamma\ kn\ iid\ np\ dp\ nfs$ says $kn$ is
declared as an \emph{eliminator constant together with its block}: its body is one of the two
canonical eliminator shapes (Definition~\ref{def:ElimBody}) and $iid$'s inductive block is
declared in $\Gamma$ with the same numbers. $\code{RuntimeKey}\ \Gamma\ kn$ is the existential
closure of \code{ElimDecl} over its four data: a key the pass \emph{consumes} rather than emits,
whose occurrences become \code{case} nodes and whose declaration the emitted environment prunes.
A constructor constant is deliberately not a runtime key, because the erasure relation already
emits the constructor node. Finally $\code{ClosedBodies}\ \Gamma$ says every declared body is de
Bruijn-closed and $\code{FVarFreeBodies}\ \Gamma$ says every declared body has no free variable;
the two are logically independent, because closedness says nothing at an \code{fvar} node.
The Lean declarations are \code{DefnDecl}, \code{IndBodyOf}, \code{ElimDecl},
\code{RuntimeKey}, \code{ClosedBodies} and \code{FVarFreeBodies}; \code{DefnDecl} is literally
an equation on \code{LBTerm.envLookup}, so it is a \emph{definitional} query and not an
inductive characterisation.
\end{definition}

\begin{remark}
\code{ClosedBodies} is the second conjunct of \code{LBWfSpec}
(Definition~\ref{def:LBWfSpec}) and the single environment-side premise the pass metatheory
carries. The \code{RuntimeKey} guard is load-bearing for \emph{inversion}, not only for
emission: \code{ElimDecl} implies \code{DefnDecl}, so without it nothing would keep a
\code{fix} image off an eliminator constant.
\end{remark}

\begin{lemma}[Eliminator data are determined by the key]
\label{lem:ElimDecl-uniq}
\lean{LeanToLambdaBox.ElimDecl.uniq}
\leanok
\uses{def:ElimDecl}
Two eliminator declarations at one key agree on all four data: if
$\code{ElimDecl}\ \Gamma\ kn\ iid\ np\ dp\ nfs$ and
$\code{ElimDecl}\ \Gamma\ kn\ iid'\ np'\ dp'\ nfs'$ then $iid = iid'$, $np = np'$, $dp = dp'$
and $nfs = nfs'$. The Lean statement is \code{ElimDecl.uniq}, whose conclusion is the
fourfold conjunction rather than an equality of the whole datum.
\begin{proof}
\leanok
\uses{def:ElimDecl, def:ElimBody, def:mkElimBody}
Environment lookup is a function, so both readings speak about one body. The two eliminator
shapes are each injective in $(iid, np, dp, nfs)$ --- the alternative list
$\code{elimAlts}\ nfs$ determines $nfs$, the telescope length determines $dp$, and the
\code{case} node carries $(iid, np)$ --- and they are never equal to one another, one being a
lambda-telescope over a \code{case} and the other a \code{fix}.
\end{proof}
\end{lemma}

\begin{definition}[Block rewriting and block closure]
\label{def:ConstToFVar}
\lean{LeanToLambdaBox.ConstToFVar, LeanToLambdaBox.CloseConstAt}
\leanok
\uses{def:LBTerm, def:Kername, def:closeFix}
$\code{ConstToFVar}\ kns\ ids$ is the congruence over \lbox{} syntax that replaces the block's
own constants by the block's fix variables: $\code{const}\ kns[j]$ maps to
$\code{fvar}\ ids[j]$ (arm \code{hit}), a constant outside the block maps to itself (arm
\code{miss}), and every other node is a congruence in which binder \emph{names} are free and
only the branch binder \emph{counts} are pinned. A \code{fix} node maps to itself without
recursing, which is sound because the source side of the pass declares no block member as a
\code{fix} and any nested \code{fix} comes from another block whose members are none of $kns$.
$\code{CloseConstAt}\ kns\ ids\ t\ u$ is the block closure of one body: $u$ is
$\code{closeFix}\ ids\ 0$ applied to a \code{ConstToFVar} image of $t$, exactly the fold the
eraser performs when it stores a block. In Lean, \code{ConstToFVar} is an inductive relation
and \code{CloseConstAt} is the existential composing it with \code{closeFix}.
\end{definition}

\begin{remark}
\code{CloseConstAt} is phrased through the existing \code{closeFix} precisely so that
Theorem~\ref{thm:closeFix_substList_fixSubst} applies verbatim and no kername-keyed twin of the
fixpoint theorems is needed.
\end{remark}

\begin{definition}[The pass relation]
\label{def:Lower}
\lean{LeanToLambdaBox.Lower, LeanToLambdaBox.LowerAlt, LeanToLambdaBox.LowerAlts}
\leanok
\uses{def:ElimDecl, def:ConstToFVar, def:LBTerm, def:LBTerm-mkApps, def:FixDef,
def:BinderName, def:GlobalDeclarations}
$\Lower\ \Gamma\ t\ t'$ says: $t'$ is a \lbox{} term the shipping eraser may emit for the
specification term $t$ over the specification environment $\Gamma$. The relation has fourteen
arms --- eleven congruence, one redex, two recursion --- and is indexed by $\Gamma$ and by
nothing else: no source expression, no typing context, no erasure run state, which is what keeps
it a statement about \lbox{} alone. The eleven congruences are the identity up to binder
\emph{names}, since the target semantics reads only their number; the \code{const} arm is
guarded by $\neg\,\code{RuntimeKey}\ \Gamma\ kn$, so an eliminator constant, whose declaration
the pass prunes, has no \code{const} image. The single redex arm \code{elimApp} turns a
saturated eliminator application, of head $kn$ applied to $\overline{pre}$, a discriminant
$disc$, the minor premises and any over-application, into
$\code{case}\ (iid, np)\ disc'\ \overline{alts}$ applied to the lowered over-application: the
$dp$ arguments before the discriminant --- parameters, motive, and for a recursor the minors'
prefix --- are \textbf{dropped}, each minor's lambda-chain is peeled into an alternative by
\code{LowerAlt}, and the branch arities are those the iota rule will read. The two recursion
arms relate a block member's \emph{constant} (\code{fixConst}, the call site) and the member's
\emph{specification body} (\code{fixBody}, the value side the delta step needs once the
environment has unfolded the constant) to the very same $\code{fix}\ defs\ j$ node.
\code{LowerAlt} is branch peeling: a minor's lambda-chain of length $nf$ becomes an
alternative's binder list of length $nf$ with the body related by the pass, names free.
\code{LowerAlts} is its pointwise lift over a block's field arities. The Lean declarations are
\code{Lower} and \code{LowerAlt}, defined mutually, and \code{LowerAlts}; the list premises are
written in indexed form (a length equation plus a bounded quantifier) because a
\code{List.Forall2} premise would be a nested inductive occurrence the kernel rejects.
\end{definition}

The arm-by-arm correspondence to MetaRocq, as the development records it, is the following.

\begin{center}
\begin{tabular}{lll}
\textbf{Arm} & \textbf{Relates} & \textbf{Counterpart} \\
\hline
\code{box} & box to box & \code{optimize} \\
\code{bvar} & an index to itself & \code{optimize} \\
\code{fvar} & a free variable to itself & \code{optimize} \\
\code{prim} & a primitive to itself & \code{optimize} \\
\code{const} & a non-runtime-key constant to itself & \code{optimize} \\
\code{lambda} & under the binder & \code{optimize} \\
\code{letIn} & value and body & \code{optimize} \\
\code{app} & head and argument & \code{optimize} \\
\code{proj} & the projected term & \code{optimize} \\
\code{construct} & the node's arguments pointwise & \code{optimize} \\
\code{case} & discriminant and branches, arities pinned & \code{iota\_red} \\
\code{elimApp} & a saturated eliminator spine to a \code{case} spine & \code{iota\_red} \\
\code{fixConst} & a block member's constant to the block's node & \code{fixSubst} \\
\code{fixBody} & a block member's body to the block's node & \code{fixSubst} \\
\end{tabular}
\end{center}

\begin{remark}
Four things a reader of [S] might expect are absent or different. (i)~The relation is
deliberately \textbf{non-deterministic} at a block member, where \code{const} and
\code{fixConst} both apply, so ``the pass emits $t'$'' must always be read as ``there is an
emission $t'$ such that''. (ii)~There is no constructor-introduction arm --- the erasure
relation already emits the nullary constructor node and the arguments arrive through
application --- and no eta-expansion arm for under-applied constructors or eliminators: the two
that once existed were refuted, and what they covered is now a decidable coverage restriction
(tag N19) together with the shipping finding F-ETA2. (iii)~There is no \code{fix} congruence arm
at all, the specification environment declaring no \code{fix}; that is what
Proposition~\ref{prop:Lower-source_const} proves. (iv)~\code{elimApp} discards the arguments
before the discriminant with no premise about them, which is sound for a \emph{forward}
simulation only: the pass is explicitly one-directional and is not a bisimulation.
\end{remark}

\begin{definition}[The block hypothesis bundle]
\label{def:LowerBlock}
\lean{LeanToLambdaBox.LowerBlock}
\leanok
\uses{def:Lower, def:ElimDecl, def:ConstToFVar, def:FixDef, def:hasFVar}
The premises the two recursion arms share, packaged as a twelve-field structure. A block
consists of $kns$ (the members' kernames, pairwise distinct), $bs$ (their specification bodies,
each declared in $\Gamma$), $bs'$ (their pass images), one \textbf{block-shared} list of fresh
fix variables $ids$ --- one list per block, as the eraser mints it, not one per member --- and
the emitted definitions $defs$, each of whose bodies is the block closure of the corresponding
member of $bs'$. Three fields are load-bearing beyond bookkeeping: \code{hfresh} (no fix
variable already occurs in a lowered body), \code{hrarg} (every emitted principal argument index
is $0$, the only value the emitter produces and the value at which one application step consumes
exactly one argument), and \code{hfl} (every emitted definition body is lambda-headed, which is
the \code{fixLambda} clause of Definition~\ref{def:LBWfPeregrine} read at this one block). The
Lean declaration is \code{LowerBlock}, a structure in \code{Prop} whose remaining nine fields are
length equations, distinctness of $kns$ and of $ids$, declaredness of each body, and the two
pointwise clauses relating $bs$ to $bs'$ and $bs'$ to the stored definition bodies.
\end{definition}

\begin{remark}
This is a \emph{hypothesis bundle}, not a proved fact: wherever the structure appears as a
hypothesis, each field is an assumed premise. It is not vacuous ---
Lemma~\ref{lem:LowerFixFixture-lowerfix_app_step} inhabits all twelve fields at a genuine
two-member mutual block. The block-shared $ids$ is forced: per-member existential names cannot
feed Theorem~\ref{thm:closeFix_substList_fixSubst}, whose freshness clause is against
\emph{every} node of the block, and closedness cannot supply it because closedness is vacuously
true at an \code{fvar} node. The \code{hrarg} field is why the emitter's own source comment,
calling the principal-argument index computationally irrelevant, is wrong under this evaluation
relation.
\end{remark}

\begin{lemma}[The three arms in packaged form]
\label{lem:Lower-fixConst-prime}
\lean{LeanToLambdaBox.Lower.fixConst', LeanToLambdaBox.Lower.fixBody',
LeanToLambdaBox.Lower.elimApp'}
\leanok
\uses{def:Lower, def:LowerBlock, def:ElimDecl}
The three arms whose premises had to be inlined, re-exposed in packaged form: \code{fixConst}
and \code{fixBody} through \code{LowerBlock}, and \code{elimApp} through \code{LowerAlts}. Every
downstream construction of a fixpoint or eliminator derivation goes through these. The Lean
names are \code{Lower.fixConst'}, \code{Lower.fixBody'} and \code{Lower.elimApp'}; the inlining
they undo is forced by the Lean kernel, which rejects a structure premise mentioning the
inductive being defined.
\begin{proof}
\leanok
\uses{def:Lower, def:LowerBlock}
Each is the corresponding constructor applied to the projections of the bundle.
\end{proof}
\end{lemma}

\section{Metatheory of the pass}
\label{sec:lowering:meta}

\begin{lemma}[A block's node is closed and variable-free]
\label{lem:lbClosed_fix_of_block}
\lean{LeanToLambdaBox.lbClosed_closeFix, LeanToLambdaBox.ConstToFVar.closed,
LeanToLambdaBox.lbClosed_fix_of_block, LeanToLambdaBox.constToFVar_not_hasFVar,
LeanToLambdaBox.closeConstAt_not_hasFVar, LeanToLambdaBox.fixNode_not_hasFVar,
LeanToLambdaBox.toBvar_fixNode}
\leanok
\uses{def:ConstToFVar, def:LBClosed, def:closeFix, def:hasFVar, def:toBvar}
The \code{fix} node a block builds is inert. Closing at base $0$ turns a closed term into one
closed at the block's width, and rewriting block constants to fix variables preserves closedness
at every bound (neither \code{const} nor \code{fvar} binds an index); hence a block's \code{fix}
node is closed at every bound once every lowered body is closed. On the free-variable side, the
rewriting introduces only the block's own identifiers, the block closure of a variable-free body
is variable-free at \emph{every} variable including the block's own, hence the block's
\code{fix} node has no free variable at all, hence abstraction is the identity on it. The Lean
statements are \code{lbClosed\_closeFix}, \code{ConstToFVar.closed},
\code{lbClosed\_fix\_of\_block}, \code{constToFVar\_not\_hasFVar},
\code{closeConstAt\_not\_hasFVar}, \code{fixNode\_not\_hasFVar} and \code{toBvar\_fixNode}; the
block-level ones take the bundle's length and closure fields as separate hypotheses rather than
the bundle itself, so that they also apply inside its construction.
\begin{proof}
\leanok
\uses{def:ConstToFVar, def:closeFix, lem:toBvar_eq_of_not_hasFVar,
lem:closeFixFold_fvar_head}
Induction over the rewriting relation for the closedness and variable clauses; the block-level
statements then read each definition body as a $\code{closeFix}\ ids\ 0$ image through the
bundle's closure field and observe that closing at base $0$ closes exactly the block's own width
and removes exactly the scheduled identifiers. The abstraction equation is the general
``abstraction is the identity where the variable does not occur'' applied to that.
\end{proof}
\end{lemma}

\begin{proposition}[The pass invents no de Bruijn index]
\label{prop:Lower-closed}
\lean{LeanToLambdaBox.Lower.closed}
\leanok
\uses{def:Lower, def:ElimDecl, def:LBClosed}
If every constant declared in $\Gamma$ has a closed body, then $\Lower\ \Gamma\ s\ t$ preserves
closedness at every bound: $\code{LBClosed}\ s\ k$ implies $\code{LBClosed}\ t\ k$. The Lean
statement \code{Lower.closed} quantifies the bound $k$ inside the conclusion, so the induction
can raise it under binders.
\begin{proof}
\leanok
\uses{def:Lower, lem:lbClosed_fix_of_block, lem:LBClosed-substList}
Induction with the mutual recursor, using a \code{LowerAlt} motive that adds the alternative's
binder count to the bound. No congruence arm invents an index; the \code{elimApp} arm splits and
reassembles a spine, and the two recursion arms read their definitions off $\Gamma$ --- which is
exactly why the environment premise is there --- and discharge through
Lemma~\ref{lem:lbClosed_fix_of_block}.
\end{proof}
\end{proposition}

\begin{proposition}[The pass commutes with shifting and substitution]
\label{prop:Lower-subst_comm}
\lean{LeanToLambdaBox.Lower.shift_comm, LeanToLambdaBox.Lower.subst_comm,
LeanToLambdaBox.Lower.substList_comm}
\leanok
\uses{def:Lower, def:ElimDecl, def:LBTerm-shift, def:LBTerm-subst}
Over an environment with closed bodies, the pass commutes with de Bruijn shifting, with
substitution of related terms, and with a whole simultaneous substitution of pointwise related
terms. These are the laws a forward simulation spends at its beta, zeta and iota steps: after
the source contracts a redex, the target's contractum is still a pass image. The Lean statements
are \code{Lower.shift\_comm}, \code{Lower.subst\_comm} and \code{Lower.substList\_comm}; the
substitution laws require the substituends themselves to be related, not merely closed.
\begin{proof}
\leanok
\uses{def:Lower, prop:Lower-closed, lem:lbClosed_fix_of_block}
Induction on the derivation. The pass reads no index, so every congruence arm is immediate; at
the two recursion arms both sides are closed by Lemma~\ref{lem:lbClosed_fix_of_block}, so the
shift and the substitution are the identity there. The \code{bvar} case of the substitution law
is where the shift law is spent. The list form is the iterated single form and is what the iota
step's branch application needs, since the iota rule substitutes a constructor's reversed field
list.
\end{proof}
\end{proposition}

\begin{proposition}[The pass commutes with binder abstraction]
\label{prop:Lower-abstract}
\lean{LeanToLambdaBox.Lower.noFVar, LeanToLambdaBox.Lower.abstract,
LeanToLambdaBox.LowerAlt.abstract}
\leanok
\uses{def:Lower, def:ElimDecl, def:toBvar, def:hasFVar}
Over an environment whose declared bodies have no free variables, the pass introduces no free
variable, and it commutes with abstraction of a free variable at any level:
$\Lower\ \Gamma\ s\ t$ gives
$\Lower\ \Gamma\ (\code{toBvar}\ x\ lvl\ s)\ (\code{toBvar}\ x\ lvl\ t)$. The alternative form
shifts the abstraction level by the branch's field count. The Lean statements are
\code{Lower.noFVar}, \code{Lower.abstract} and \code{LowerAlt.abstract}; they carry
\code{FVarFreeBodies} and \emph{not} \code{ClosedBodies}, and the next lemma shows the
difference is real.
\begin{proof}
\leanok
\uses{def:Lower, lem:lbClosed_fix_of_block}
The variable-freedom statement comes first and independently: both recursion arms discharge
through Lemma~\ref{lem:lbClosed_fix_of_block} without using the induction hypothesis on the
block. The abstraction law then carries no cutoff bound and no closedness premise, because
abstraction never reads or rewrites an existing index; the two recursion arms do not recurse but
rewrite by the equation saying the block's node is inert under abstraction, whose input is the
variable-freedom statement at the block's own bodies. The alternative form is obtained by
induction on the field count rather than by re-running the fourteen arms.
\end{proof}
\end{proposition}

\begin{remark}
The abstraction law is needed because the eraser closes every binder with the same abstraction
operator, and the erasure relation closes its image with it too, so the pass factor must follow.
\code{LowerAlt.abstract} has no consumer elsewhere in the tree.
\end{remark}

\begin{lemma}[Closedness does not replace variable-freedom]
\label{lem:FVarFixture-noFVar_needs_fvarFree}
\lean{LeanToLambdaBox.FVarFixture.noFVar_needs_fvarFree,
LeanToLambdaBox.FVarFixture.abstract_needs_fvarFree,
LeanToLambdaBox.FVarFixture.closed_cxEnv, LeanToLambdaBox.FVarFixture.not_fvarFree,
LeanToLambdaBox.FVarFixture.blk, LeanToLambdaBox.FVarFixture.lower_const}
\leanok
\uses{def:Lower, def:ElimDecl, def:hasFVar}
Proposition~\ref{prop:Lower-abstract} is \textbf{false} if the variable-freedom premise is
weakened to closedness. Both refutations run on one one-member block whose \emph{declared} body
is a lambda over a stray free variable: the environment satisfies \code{ClosedBodies} ---
closedness says nothing at an \code{fvar} node --- but not \code{FVarFreeBodies}, and the
\code{fixConst} arm relates the variable-free \code{const} to a \code{fix} node carrying the
declared body's stray variable. The Lean statements
\code{FVarFixture.noFVar\_needs\_fvarFree} and \code{FVarFixture.abstract\_needs\_fvarFree} are
negations of the universally quantified weakened laws, so they refute the general statement and
not merely one instantiation.
\begin{proof}
\leanok
\uses{def:LowerBlock, lem:Lower-fixConst-prime, prop:Lower-closed}
Exhibit the one-entry environment and the one-member block; the \code{fixConst} arm fires at it,
which refutes the variable statement directly. For the abstraction statement, abstracting the
stray variable under the block's own binders would manufacture an out-of-range de Bruijn index,
and since the pass preserves closedness no arm can relate the abstracted pair.
\end{proof}
\end{lemma}

\section{The inversion kit}
\label{sec:lowering:inversion}

\begin{lemma}[Target-side inversion and spine congruence]
\label{lem:Lower-target_fix}
\lean{LeanToLambdaBox.Lower.target_box, LeanToLambdaBox.Lower.target_bvar,
LeanToLambdaBox.Lower.target_fvar, LeanToLambdaBox.Lower.target_prim,
LeanToLambdaBox.Lower.target_const, LeanToLambdaBox.Lower.target_fix,
LeanToLambdaBox.Lower.target_construct, LeanToLambdaBox.Lower.mkApps}
\leanok
\uses{def:Lower, def:LowerBlock, def:ElimDecl, def:LBTerm-mkApps}
What a target of a given shape can have come from. An atom --- box, index, free variable,
primitive --- comes from the same atom: no arm introduces or erases a box, renumbers an index,
renames a free variable or rewrites a primitive. A constant in the target comes from the same
constant, and that constant is not a runtime key. A constructor node comes from the same node,
argument by argument. A \code{fix} node comes from one of the two recursion arms and carries a
whole \code{LowerBlock}: either the member's constant or the member's specification body.
Conversely the relation lifts from a head and its arguments to a whole application spine. The
Lean statements are the \code{Lower.target\_*} family and \code{Lower.mkApps}; each inversion is
phrased with the target shape as a separate equation hypothesis, which is what lets it be
applied to a derivation whose target is only later known.
\begin{proof}
\leanok
\uses{def:Lower, def:LowerBlock}
Each is a case analysis on the derivation with the target shape fixed; the arms whose targets
are a \code{case} spine or a \code{fix} node are excluded by the shape, and the surviving arm is
the congruence of the same node. The \code{fix} case keeps the whole bundle, which is what the
delta and beta steps through a block run on. The spine congruence is an induction on the
argument list through the \code{app} arm.
\end{proof}
\end{lemma}

\begin{proposition}[A block member's specification body is a lambda]
\label{prop:LowerBlock-lambda_of_fixLambda}
\lean{LeanToLambdaBox.LowerBlock.lambda_of_fixLambda,
LeanToLambdaBox.LowerBlock.targetLambda_of_fixLambda, LeanToLambdaBox.Lower.source_isLambda,
LeanToLambdaBox.isLambda_toBvar, LeanToLambdaBox.isLambda_closeFix,
LeanToLambdaBox.ConstToFVar.isLambda_eq}
\leanok
\uses{def:LowerBlock, def:Lower, def:atomValue}
Unconditionally --- with no environment-side premise --- a block member's specification body is
lambda-headed. Abstraction, the block-closing fold and the constant-to-fixvar rewriting are each
faithful on the head constructor (a lambda stays a lambda and a non-lambda stays a non-lambda),
so the \code{hfl} field reads back first to the lowered bodies and then, through the pass itself,
to the specification bodies; the pass step is legitimate because a lambda-headed target comes
from a lambda-headed source, \code{lambda} being the only arm whose target is a lambda. The Lean
statements are \code{LowerBlock.lambda\_of\_fixLambda},
\code{LowerBlock.targetLambda\_of\_fixLambda} and \code{Lower.source\_isLambda}, together with
the three head-faithfulness equations; all are stated through the Boolean head test
\code{isLambda} rather than an existential.
\begin{proof}
\leanok
\uses{def:LowerBlock, def:Lower, def:ConstToFVar, def:closeFix, def:toBvar}
The head-faithfulness of the three syntactic operations is a one-line case split on the head
constructor in each case. \code{Lower.source\_isLambda} is a case analysis on the derivation:
\code{elimApp}'s target is a \code{case} spine and the two recursion arms' target is a
\code{fix}, so \code{lambda} is the only surviving arm. Composing the three transports with the
bundle's \code{hfl} field gives the statement.
\end{proof}
\end{proposition}

\begin{remark}
This single fact is what lets the entire source-side inversion kit dispense with an environment
guard, since it rules a \code{fix} image out at every source shape but a lambda.
\end{remark}

\begin{lemma}[Only a constant or a lambda has a block image]
\label{lem:Lower-ne_fix_of_block}
\lean{LeanToLambdaBox.Lower.notFix_of_block, LeanToLambdaBox.Lower.ne_fix_of_block}
\leanok
\uses{def:Lower, def:atomValue}
If $\Lower\ \Gamma\ s\ (\code{fix}\ defs\ j)$ then $s$ is a constant or is lambda-headed;
contrapositively, a source that is neither has no \code{fix} image. Both Lean statements,
\code{Lower.notFix\_of\_block} and \code{Lower.ne\_fix\_of\_block}, are premise-free in the
environment.
\begin{proof}
\leanok
\uses{lem:Lower-target_fix, prop:LowerBlock-lambda_of_fixLambda}
Invert the target through Lemma~\ref{lem:Lower-target_fix}; the \code{fixConst} reading gives a
constant, and the \code{fixBody} reading gives a declared block-member body, which
Proposition~\ref{prop:LowerBlock-lambda_of_fixLambda} pins to a lambda.
\end{proof}
\end{lemma}

\begin{proposition}[Source-side inversion]
\label{prop:Lower-source_const}
\lean{LeanToLambdaBox.Lower.source_box, LeanToLambdaBox.Lower.source_bvar,
LeanToLambdaBox.Lower.source_fvar, LeanToLambdaBox.Lower.source_prim,
LeanToLambdaBox.Lower.source_letIn, LeanToLambdaBox.Lower.source_proj,
LeanToLambdaBox.Lower.source_construct, LeanToLambdaBox.Lower.source_case,
LeanToLambdaBox.Lower.source_fix, LeanToLambdaBox.Lower.source_lambda,
LeanToLambdaBox.Lower.source_const, LeanToLambdaBox.Lower.source_construct_nil,
LeanToLambdaBox.LowerAlt.arity}
\leanok
\uses{def:Lower, def:ElimDecl}
For each source node shape, the possible targets, with \textbf{no} environment-side guard. An
atom goes to itself; a \code{let} to a \code{let} with related value and body; a projection to a
projection with the same triple; a constructor node argument by argument; a \code{case} to a
\code{case} with the same inductive, parameter count, branch count and branch arities. A
\code{fix} in the source has \emph{no image at all} --- the statement concludes falsity ---
because the relation has no \code{fix} congruence arm, \code{fixConst}'s source is a constant
and \code{fixBody}'s source is pinned to a lambda by the block. A lambda goes to a lambda
\emph{or} to a block's \code{fix} node. A constant goes to itself \emph{or} to a block's
\code{fix} node, and in \emph{both} cases it is not a runtime key, which is precisely what lets
an eliminator declaration at that key refute both images. A branch peeling pins the
alternative's binder count to the field arity and nothing else --- exactly what the iota rule
reads. The Lean statements are the \code{Lower.source\_*} family together with
\code{LowerAlt.arity}, each phrased with the source shape as a separate equation hypothesis.
\begin{proof}
\leanok
\uses{def:Lower, lem:Lower-ne_fix_of_block, prop:LowerBlock-lambda_of_fixLambda}
Case analysis on the derivation with the source shape fixed. At every shape but a constant and a
lambda the two recursion arms are excluded premise-free by
Lemma~\ref{lem:Lower-ne_fix_of_block}, and the \code{elimApp} arm is excluded because its source
is a spine headed by a constant. At a constant both readings survive, and the \code{fixBody}
reading is excluded by the block's lambda-headedness (it would force a constant to be
lambda-headed), leaving the runtime-key guard common to both.
\end{proof}
\end{proposition}

\begin{remark}
The application case is not here: \code{Lower.source\_app} lives in
\code{ErasesCorrect/Steps.lean} and is owned by Chapter~\ref{chap:correctness}
(Lemma~\ref{lem:Lower-source_app}).
\end{remark}

\begin{lemma}[The eliminator arm fires]
\label{lem:LowerElimFixture-elimApp_fires}
\lean{LeanToLambdaBox.LowerElimFixture.elimDecl,
LeanToLambdaBox.LowerElimFixture.runtimeKey_elimKn,
LeanToLambdaBox.LowerElimFixture.elimApp_fires}
\leanok
\uses{def:Lower, def:ElimDecl}
On a hand-built two-constructor inductive, with constructors of $0$ and $1$ fields, and its
\code{casesOn} constant, both conjuncts of the eliminator declaration hold; hence that constant
\emph{is} a runtime key, so neither the \code{const} arm nor the \code{fixConst} arm relates it;
and the \code{elimApp} arm fires at it with a non-empty over-application, producing a
\code{case} node applied to one extra argument. The three Lean statements are stated at a fully
explicit environment, so each is a closed fact about concrete data.
\begin{proof}
\leanok
\uses{def:ElimDecl, def:mkElimBody, lem:Lower-fixConst-prime}
The environment is given explicitly, so both halves of the eliminator declaration are
definitional; the runtime-key statement is the existential over them; the arm is applied with
the branch peelings discharged by the two \code{LowerAlt} constructors.
\end{proof}
\end{lemma}

\begin{lemma}[Lambda-headedness must be keyed on the block]
\label{lem:LowerCtorBodyFixture-lowerBlock_needs_lambda_bodies}
\lean{LeanToLambdaBox.LowerCtorBodyFixture.lowerBlock_needs_lambda_bodies}
\leanok
\uses{def:LowerBlock}
A definition whose body is a nullary constructor node --- the shape of \code{Unit.unit}, which
every one of the five benchmark programs declares --- is \textbf{not} a block of the pass: the
lambda-headedness field fails, because the emitted body is the constructor node, not a lambda.
The Lean statement is the negation of a \code{LowerBlock} claim at that concrete one-member
block.
\begin{proof}
\leanok
\uses{def:LowerBlock}
Project the \code{hfl} field at index $0$ and evaluate: the stored body's head is a constructor
node.
\end{proof}
\end{lemma}

\begin{remark}
This is why lambda-headedness is keyed on the blocks a run actually builds rather than
quantified over every block over the environment --- an environment-wide condition would be
refuted by one such declaration. It is recorded as coverage restriction N22.
\end{remark}

\section{The fixpoint round trip}
\label{sec:lowering:fix}

\begin{definition}[Substituting for free variables; the static fix substitution]
\label{def:substFVar}
\lean{LeanToLambdaBox.substFVar, LeanToLambdaBox.substFVarList, LeanToLambdaBox.substFix}
\leanok
\uses{def:LBTerm, def:FixDef, def:LBTerm-fixSubst}
$\code{substFVar}\ x\ s$ replaces the free variable $x$ by $s$ throughout, with no de Bruijn
bookkeeping under binders --- every law about it carries closedness of $s$, and a closed term is
fixed by shifting. \code{substFVarList} iterates it, innermost binding first.
$\code{substFix}\ ids\ defs$ substitutes the block's own nodes for the block's fix variables,
sending $\code{fvar}\ ids[j]$ to $\code{fix}\ defs\ j$; it is the \emph{static} counterpart of
the dynamic substitution the target's guarded fixpoint rule performs. The Lean definitions are
\code{substFVar}, \code{substFVarList} and \code{substFix}, the last built by zipping $ids$ with
its own index.
\end{definition}

\begin{remark}
The same file proves closedness preservation for both substitution forms and the ``identity
where the variable does not occur'' laws, which the transports below consume; they are named in
the proofs that spend them rather than given nodes of their own.
\end{remark}

\begin{lemma}[Abstraction against substitution: the two primitives]
\label{lem:subst_toBvar_self}
\lean{LeanToLambdaBox.subst_toBvar_self, LeanToLambdaBox.subst_toBvar_succ,
LeanToLambdaBox.substList_toBvar, LeanToLambdaBox.substList_toBvar'}
\leanok
\uses{def:substFVar, def:toBvar, def:LBClosed, def:LBTerm-subst, def:hasFVar}
Two cases, by how the abstraction level compares with the substitution depth. When they agree,
the freshly created binder is consumed: abstracting $x$ to level $d$ and then substituting a
closed $s$ at depth $d$ is the same as substituting $s$ for $x$ directly, provided $t$ has no
loose index at $d$ or above. When the abstraction level is above the depth the binder survives,
one lower --- provided $s$ is closed and does not mention $x$, otherwise the abstraction on the
right would reach into the substituted term. The simultaneous form pushes a substitution of
closed, $x$-free terms past an abstraction, dropping the level by exactly the number of
substitutions. The Lean statements are \code{subst\_toBvar\_self}, \code{subst\_toBvar\_succ},
\code{substList\_toBvar} and its primed variant; the level and the depth are quantified inside
the statement so the binder cases can raise them.
\begin{proof}
\leanok
\uses{def:substFVar}
Structural induction on the term, the level quantified inside the statement so the binder cases
can hand the incremented level to the induction hypothesis. In the level-agreeing case the only
index at depth $d$ is the one the abstraction just created, which is the closedness premise's
job, and the substitution's internal shift is the identity because the substituted term is
closed. The simultaneous form is the iterated step case.
\end{proof}
\end{lemma}

\begin{lemma}[What closing removes, and what it schedules]
\label{lem:hasFVar_toBvar}
\lean{LeanToLambdaBox.hasFVar_toBvar, LeanToLambdaBox.not_hasFVar_closeFixFold,
LeanToLambdaBox.not_hasFVar_closeFix, LeanToLambdaBox.closeFixFold_append,
LeanToLambdaBox.closeFix_cons, LeanToLambdaBox.substFVarList_zipIdx_fvar,
LeanToLambdaBox.substFix_fvar_getElem}
\leanok
\uses{def:substFVar, def:toBvar, def:closeFix, def:hasFVar}
Abstraction removes exactly one variable and creates none: if $x$ occurs in
$\code{toBvar}\ y\ lvl\ t$ then $x \neq y$ and $x$ already occurred in $t$. Hence the closing
fold closes everything it is given --- a term whose free variables all lie in $ids$ closes to a
term with no free variable at all. The fold also has the structural recursion its definition
hides: peeling one identifier needs an append law for the fold first, and the head identifier is
the outermost abstraction. Dually, the static fix substitution does what it is meant to: it
sends $ids[j]$ to the block's $j$-th node, under distinctness of $ids$ and the block's own
variable-freedom. The Lean statement \code{substFVarList\_zipIdx\_fvar} is generalised over the
window's base index, which is what makes the induction go through.
\begin{proof}
\leanok
\uses{def:substFVar, def:closeFix}
The abstraction fact is a structural induction: the only node abstraction rewrites is an
\code{fvar} matching $y$. The fold statements follow by induction over the identifier list using
the append law. The substitution computation is generalised over the window's base index;
variable-freedom of the block is what lets an already-substituted sibling pass through the
remaining substitutions untouched.
\end{proof}
\end{lemma}

\begin{remark}
\code{not\_hasFVar\_closeFix} is what makes the stored block variable-free \emph{without
assuming it}, since the block's opened bodies mention only the run's own fix variables and the
emitter abstracts exactly those.
\end{remark}

\begin{theorem}[Unfolding a closed block undoes the closing]
\label{thm:closeFix_substList_fixSubst}
\lean{LeanToLambdaBox.closeFix_substList_fixSubst,
LeanToLambdaBox.closeFix_substList_fixSubst_gen,
LeanToLambdaBox.closeFix_substList_fixSubst_fires,
LeanToLambdaBox.closeFix_substList_fixSubst_fires_value}
\leanok
\uses{def:substFVar, def:closeFix, def:LBTerm-fixSubst, def:LBTerm-subst, def:LBClosed,
def:hasFVar}
Substituting a block into a body the emitter closed with $\code{closeFix}\ ids\ 0$ --- which is
exactly what the target's guarded fixpoint rule performs --- is the same as substituting
$\code{fix}\ defs\ j$ for $ids[j]$ in the \emph{open} body. The hypotheses are a length match,
that every node of the block is closed, that the block does not itself mention the fix variables
(automatic when the block was produced by the closing fold), and that the body is closed.
Notably, distinctness of $ids$ is \textbf{not} needed and is not assumed. The Lean statement is
\code{closeFix\_substList\_fixSubst}, proved from a form generalised over the window's base
index, and the two \code{\_fires} theorems discharge the hypothesis class on a concrete
self-looping block.
\begin{proof}
\leanok
\uses{lem:subst_toBvar_self, lem:hasFVar_toBvar, def:substFVar}
The general form peels one identifier at a time, the surviving suffix of the block starting at
the window's base; each peel is the cons law for the fold followed by the simultaneous
abstraction-against-substitution law, whose two premises --- closedness and $x$-freedom of every
substituted node --- are the two hypotheses. Two non-vacuity theorems close the argument: the
premises are discharged on a concrete self-looping block, and the common value is the block
itself rather than a vacuous index.
\end{proof}
\end{theorem}

\begin{remark}
The module docstring records why the variable-freedom hypothesis is not slack: with two
identifiers and one occurring free in the block, the two sides genuinely differ. This is the
single equation that lets the delta arm of a forward simulation turn a target fixpoint unfolding
into an ordinary statement about the source body.
\end{remark}

\begin{proposition}[A block's dynamic unfolding inverts its static closure]
\label{prop:LowerBlock-substList_fixSubst}
\lean{LeanToLambdaBox.LowerBlock.substList_fixSubst, LeanToLambdaBox.LowerBlock.not_hasFVar_fix,
LeanToLambdaBox.LowerBlock.lbClosed_fix, LeanToLambdaBox.not_hasFVar_closeFix_of_mem}
\leanok
\uses{def:LowerBlock, def:substFVar, def:ElimDecl, def:closeFix}
Theorem~\ref{thm:closeFix_substList_fixSubst}, with all three of its block-side hypotheses read
off a \code{LowerBlock}: the length match from the bundle's two length fields, the closedness
from the bundle plus closedness of the declared bodies, and the freshness from the fact that
every definition body is a closing-fold image and the fold removes exactly the identifiers it
schedules. Only closedness needs anything beyond the structure, in the shape of the
\code{ClosedBodies} premise the Lean statement carries.
\begin{proof}
\leanok
\uses{thm:closeFix_substList_fixSubst, lem:hasFVar_toBvar, lem:lbClosed_fix_of_block,
prop:Lower-closed}
Variable-freedom of the stored block is the closing fold's removal property applied at each
definition body. Closedness is the block-node closedness lemma fed with closedness of the
lowered bodies, which is Proposition~\ref{prop:Lower-closed} at the declared bodies --- the
bundle carries no closedness field of its own, so this is where the environment premise enters
the block story. The conclusion is then the general theorem.
\end{proof}
\end{proposition}

\begin{theorem}[Installing the block keeps the relation]
\label{thm:Lower-constToFix}
\lean{LeanToLambdaBox.Lower.constToFix, LeanToLambdaBox.ConstToFVar.mkApps_inv}
\leanok
\uses{def:Lower, def:LowerBlock, def:ConstToFVar, def:substFVar, def:ElimDecl}
Suppose the target term $t$ is a pass image of $s$ and $t'$ is its block rewriting (siblings'
constants replaced by fix variables). Then substituting the block's own nodes back for those
variables --- exactly what a fixpoint unfolding does, once
Theorem~\ref{thm:closeFix_substList_fixSubst} has rewritten the dynamic substitution into the
static one --- lands back inside the pass at the \emph{same} source term $s$. The Lean statement
\code{Lower.constToFix} takes as its freshness premise that the fix variables do not already
occur in $t$, which is free at the call site, where it is the bundle's \code{hfresh} field.
\begin{proof}
\leanok
\uses{def:Lower, def:LowerBlock, lem:Lower-fixConst-prime,
prop:LowerBlock-substList_fixSubst, lem:hasFVar_toBvar, lem:lbClosed_fix_of_block,
def:substFVar}
Induction over the pass derivation with a branch-peeling motive. Every structural arm rewrites
by the corresponding push-through equation for the static substitution, which leaves binder
names and projection triples alone; at a rewriting \code{hit} the substitution turns the fix
variable back into the block's node and the arm's conclusion is replaced by the packaged
call-site arm; at a source already related by a recursion arm the substitution is the identity,
by the block's own variable-freedom. The application-spine arm needs the inversion of the
rewriting through a spine.
\end{proof}
\end{theorem}

\begin{remark}
The measured axiom footprint of \code{Lower.constToFix} is \code{propext} and \code{Quot.sound}
only --- no \code{sorryAx} --- and the row is a committed fixture in
\code{test/ledger.expected}. In the shipping simulation the transport is spent through
\code{Lower.appReady} (Lemma~\ref{lem:Lower-appReady}), owned by
Chapter~\ref{chap:correctness}.
\end{remark}

\begin{lemma}[The design document's form of the transport is false]
\label{lem:LowerFixFixture-constToFix_needs_freshness}
\lean{LeanToLambdaBox.LowerFixFixture.constToFix_needs_freshness}
\leanok
\uses{thm:Lower-constToFix, def:LowerBlock, def:ConstToFVar}
Theorem~\ref{thm:Lower-constToFix} with the freshness premise replaced by mere closedness of $t$
is \textbf{refutable}. At $s = t = t' = \code{fvar}\ ids[0]$ every premise of that weaker
statement holds --- closedness says nothing about free variables --- while the conclusion
demands that a free variable be related to the block's node, which no arm derives: the only
value-side recursion arm needs a declared specification body, and no declared body in the
fixture is a free variable. The Lean statement negates the universally quantified weakened form.
\begin{proof}
\leanok
\uses{lem:LowerFixFixture-lowerfix_app_step, prop:Lower-source_const}
Instantiate at the two-member fixture block and at the free variable; the conclusion is inverted
by the source-side kit at an \code{fvar}, which admits only the same \code{fvar} as image.
\end{proof}
\end{lemma}

\begin{proposition}[A member's unfolded definition is a lambda still related to its body]
\label{prop:Lower-fixUnfold}
\lean{LeanToLambdaBox.Lower.fixUnfold}
\leanok
\uses{def:Lower, def:LowerBlock, def:ElimDecl, def:LBTerm-fixSubst}
For a block and an index in range, the unfolding of the $j$-th definition's body is a lambda,
and that lambda is still related by the pass to the $j$-th member's \emph{specification} body.
The Lean statement \code{Lower.fixUnfold} needs no premise beyond the block and closedness of
the declared bodies, and its conclusion is an existential exhibiting the binder name and the
lambda's body.
\begin{proof}
\leanok
\uses{prop:LowerBlock-substList_fixSubst, thm:Lower-constToFix, prop:Lower-closed,
def:ConstToFVar}
The bundle writes the definition body as a closing-fold image, so
Proposition~\ref{prop:LowerBlock-substList_fixSubst} rewrites the dynamic substitution into the
static one; unfolding does not change the head constructor, so the lambda-headedness field gives
the result its lambda head; and Theorem~\ref{thm:Lower-constToFix} carries the relation across
the substitution.
\end{proof}
\end{proposition}

\begin{remark}
This is what the target does where the source delta-steps into a recursive body. In the tree as
it stands the theorem has no consumer: the transport the simulation's beta and delta arms
actually spend is \code{Lower.appReady}, which calls Theorem~\ref{thm:Lower-constToFix}
directly.
\end{remark}

\begin{definition}[Block lowering at the level of declarations]
\label{def:LowerFix}
\lean{LeanToLambdaBox.LowerFix}
\leanok
\uses{def:LowerBlock}
$\code{LowerFix}\ \Gamma\ kns\ bs\ defs$ says $defs$ is a lowered image of the specification
bodies $bs$ under the names $kns$, with the lowered bodies and the fix variables existentially
quantified. It tolerates an unused fixpoint binder, because the eraser decides recursiveness on
the \emph{source} body and erasure can remove the only self-reference. In Lean it is literally
the existential closure of \code{LowerBlock} over its two hidden lists, so a \code{LowerFix}
premise determines neither.
\end{definition}

\begin{remark}
\code{LowerFix} is the right disjunct of the definitional clause of $\LowerEnv$
(Definition~\ref{def:LowerEnv}), and of the corresponding clause in the cold-start shape.
\end{remark}

\begin{definition}[The three-factor composite for a sub-run inside a block]
\label{def:ErasesLBFix}
\lean{LeanToLambdaBox.ErasesLBFix}
\leanok
\uses{def:Erases, def:Lower, def:ConstToFVar, def:lean4lean-grounding}
Erasure, lowering and the block's constant-to-fixvar rewriting composed into one relation.
Inside a mutual block the eraser emits $\code{fvar}\ ids[j]$ where the source has
$\code{const}\ kns[j]$ --- a pair neither the erasure relation nor the pass relates on its own
--- so the two-factor composite cannot state what a sub-run inside the block branch concludes.
The Lean definition \code{ErasesLBFix} threads the lean4lean environment and local context of
the erasure factor. Among the three pass files, lean4lean is otherwise touched only by the
\code{LowerFixFixture} built for Lemma~\ref{lem:LowerFixFixture-lowerfix_app_step}, which
constructs a concrete lean4lean environment and whose \code{lowerfix\_erasesLBFix} proof
invokes lean4lean's typing judgement.
\end{definition}

\begin{lemma}[The fixpoint wall at a real mutual block]
\label{lem:LowerFixFixture-lowerfix_app_step}
\lean{LeanToLambdaBox.LowerFixFixture.lowerfix_nv,
LeanToLambdaBox.LowerFixFixture.lowerfix_fixConst_fires,
LeanToLambdaBox.LowerFixFixture.lowerfix_app_step,
LeanToLambdaBox.LowerFixFixture.lowerfix_erasesLBFix,
LeanToLambdaBox.LowerFixFixture.lowerfix_delta,
LeanToLambdaBox.LowerFixFixture.lowerfix_isLambda}
\leanok
\uses{def:LowerBlock, def:Lower, def:ErasesLBFix, def:WcbvEval, def:eraseFlags}
A hand-built two-member mutual block --- member $0$ is a lambda returning the box, member $1$ is
a lambda calling member $0$ on its argument --- inhabits all twelve fields of the bundle,
including the block-shared identifiers and the principal-argument field, with a real call from
member $1$ into member $0$. The call-site arm fires on it; the emitted environment delta-steps
the member's constant to the block; and one application step \emph{through the block} has the
exact shape of a pass-correctness statement: whatever the source application evaluates to, the
target application evaluates to a pass image of it. Separately, the three-factor composite is
inhabited at exactly the pair no two-factor composite relates --- the source constant against
the fix variable the eraser emits for it inside the block. All six Lean statements are at the
fixture's explicit environments and terms.
\begin{proof}
\leanok
\uses{def:LowerBlock, lem:Lower-fixConst-prime, prop:LowerBlock-substList_fixSubst,
def:Erases, thm:eval_deterministic}
The twelve fields are discharged on explicit data. The source step evaluates the application
through member $0$ to a box; the matching target run is \emph{two} guarded-fixpoint links plus
the beta steps, and determinism of evaluation identifies the source value. The composite
statement builds an erasure derivation at the source constant and the rewriting's \code{hit} arm
at the target.
\end{proof}
\end{lemma}

\begin{remark}
This is the cluster's guard against a vacuously true fixpoint story; it is the only theorem here
whose proof touches lean4lean's typing judgement, through the erasure derivation's constant
rule, and its measured axiom footprint is nevertheless free of \code{sorryAx}.
\end{remark}

\begin{lemma}[Box-freedom does not transport along the pass]
\label{lem:noBox_lower_needs_noFix}
\lean{LeanToLambdaBox.noBox_lower_needs_noFix}
\leanok
\uses{def:Lower, def:NoBox}
The naive transport of box-freedom along the pass is \textbf{false}: the call-site arm relates
the box-free \code{const} to the block's \code{fix} node, whose definitions carry the boxes of
the members' bodies. Any transport of box-freedom therefore needs a premise excluding a
\code{fix} in the target. The Lean statement negates the universally quantified transport.
\begin{proof}
\leanok
\uses{lem:LowerFixFixture-lowerfix_app_step, lem:Lower-fixConst-prime, def:NoBox}
Instantiate at the fixture block, whose member $0$ returns the box; the call-site arm relates
the box-free constant to the block, and the block's box-freedom would project to the member's
body, which is a box.
\end{proof}
\end{lemma}

\begin{remark}
This is exactly why the capstone proves box-freedom of the lowered first-order value along the
constructor-tree shape \code{FOSpine} (Definition~\ref{def:FOSpine} and
Lemma~\ref{lem:noBox_lower_of_foSpine}, Chapter~\ref{chap:correctness}) rather than by
transporting box-freedom along the pass.
\end{remark}

\section{Eliminator bodies in the specification environment}
\label{sec:lowering:elimbody}

\begin{definition}[The canonical eliminator body]
\label{def:mkElimBody}
\lean{LeanToLambdaBox.fieldArgs, LeanToLambdaBox.fieldArgs_eq, LeanToLambdaBox.elimAlts,
LeanToLambdaBox.mkElimBody, LeanToLambdaBox.mkElimBodyRec, LeanToLambdaBox.bvarsDesc}
\leanok
\uses{def:LBTerm, def:LBTerm-mkApps, def:mkLambdas, def:InductiveId, def:FixDef}
The \lbox{} body a specification environment gives an \emph{eliminator} constant.
$\code{fieldArgs}\ m$ is the list of $m$ field indices an alternative applies to its minor
premise, outermost first; $\code{elimAlts}\ nfs$ builds the alternatives of the eliminator's
\code{case} node, alternative $k$ binding its $nfs[k]$ fields and applying the $k$-th minor
premise --- which sits $nfs[k] + (|nfs| - 1 - k)$ binders out --- to them in order.
$\code{mkElimBody}\ iid\ np\ dp\ nfs$ is then the \code{casesOn} calling convention: $dp$
dropped arguments (parameters and motive), then the discriminant, then one minor premise per
constructor, all bound by an anonymous lambda-telescope, with the body a \code{case} node on the
discriminant. \code{mkElimBodyRec} is that dispatch wrapped in a one-member guarded \code{fix}
whose principal argument index is the discriminant's position. The Lean declarations are
\code{fieldArgs}, \code{elimAlts}, \code{mkElimBody} and \code{mkElimBodyRec}, with
\code{bvarsDesc} the same index run written non-recursively and \code{fieldArgs\_eq} the bridge
between the two spellings.
\end{definition}

\begin{remark}
The recursive shape's fixpoint variable is \emph{provably} unused: its body is closed, and which
fields carry a recursive call is not a function of the field counts alone, so no recursive call
can be written at this signature. The shape is retained because the predicate below needs a
right-hand side and two inversion lemmas destruct it; nothing recursive is expressed by it. Note
also that this principal-argument index is $dp$, whereas every block the pass builds has index
$0$: they are different objects, and a claim that every fixpoint definition in this development
has principal argument index $0$ would be wrong.
\end{remark}

\begin{definition}[The eliminator shape predicate]
\label{def:ElimBody}
\lean{LeanToLambdaBox.ElimBody}
\leanok
\uses{def:mkElimBody}
The \lbox{} body of an eliminator constant as a purely \emph{syntactic} shape: two constructors,
one for each canonical body, and no semantic side condition. There is deliberately no
propositional-singleton shape, because the emitted environment marks every inductive
non-propositional, so a box-discriminant \code{case} is stuck at every flag point. The Lean
declaration \code{ElimBody} is an inductive family indexed by the four data and the body.
\end{definition}

\begin{lemma}[Both eliminator shapes are closed, at real instances]
\label{lem:ElimBody-closed}
\lean{LeanToLambdaBox.ElimBody.closed, LeanToLambdaBox.fieldArgs_closed,
LeanToLambdaBox.elimAlts_closed, LeanToLambdaBox.mkElimBody_closed,
LeanToLambdaBox.mkElimBodyRec_closed, LeanToLambdaBox.natCasesOn_elimBody,
LeanToLambdaBox.boolCasesOn_elimBody, LeanToLambdaBox.decCasesOn_elimBody}
\leanok
\uses{def:ElimBody, def:mkElimBody, def:LBClosed}
Every field index is below the telescope it is read under, the alternatives are closed under the
minor-premise binders, and hence both eliminator bodies are closed terms --- which is what the
pass's eliminator arms need, with no environment-side premise at all. The construction is
exhibited at three real Lean declarations at the numbers their declarations actually have:
\code{Nat.casesOn} (no parameters, one dropped argument for the motive, minors of $0$ and $1$
fields), \code{Bool.casesOn} (two field-free minors) and \code{Decidable.casesOn} (one
parameter, so two dropped arguments, and two one-field minors). The three Lean instances are
stated at the \code{cases} constructor, so they are definitional checks rather than proofs.
\begin{proof}
\leanok
\uses{def:mkElimBody, def:mkLambdas, def:LBClosed}
Induction over the field-count list for the alternatives, bounding the largest minor index by
the telescope length; the two body statements are then the telescope closure of that, and the
shape predicate's two constructors dispatch to them. The three instances are each the
\code{cases} constructor, definitionally.
\end{proof}
\end{lemma}

\begin{remark}
\code{ElimBody.lean} carries no evaluation theory at all, because the specification environment
is never evaluated; the earlier evaluation theory was deliberately deleted. The three instances
above are the repository's ``checked instances'' criterion for this file.
\end{remark}

\section{The pass at the level of environments}
\label{sec:lowering:env}

The environment half of the pass is stated as a structure rather than as a relation on lists,
because what the target needs of the emitted environment is not a pointwise image but a
conjunction of clauses of different shapes: which bodies survive the pruning, which are pass
images, which are members of a lowered block, that inductive declarations are carried over
unchanged so that the target's arity and propositionality queries read the specification's
numbers, and that pruning only removes. That structure is $\LowerEnv$
(Definition~\ref{def:LowerEnv}), owned by Chapter~\ref{chap:erases} together with its
satisfiability witness (Lemma~\ref{lem:lowerEnv_idEnv}); its definitional clause is where
Definition~\ref{def:Lower} and Definition~\ref{def:LowerFix} appear, and its closedness clauses
are where the environment premise of this chapter's metatheory is supplied. The lambda-headedness
that the source-side inversion kit needs is deliberately \emph{not} a clause of it: it is the
block bundle's own field, keyed on the blocks the pass builds, for the reason
Lemma~\ref{lem:LowerCtorBodyFixture-lowerBlock_needs_lambda_bodies} gives. $\LowerEnv$ is one of
the two fields of the capstone's environment-bridge assumption, and is therefore \emph{assumed}
there rather than proved for a run of the shipping eraser; see Chapter~\ref{chap:trust}.

\section{Output shapes and box-freedom}
\label{sec:lowering:shape}

\begin{definition}[Box-freedom]
\label{def:NoBox}
\lean{LeanToLambdaBox.NoBox, LeanToLambdaBox.NoBoxArgs, LeanToLambdaBox.NoBoxAlts,
LeanToLambdaBox.NoBoxDefs, LeanToLambdaBox.NoBoxArgs_iff, LeanToLambdaBox.NoBoxAlts_iff,
LeanToLambdaBox.NoBoxDefs_iff, LeanToLambdaBox.NoBox_mkApps, LeanToLambdaBox.NoBox_shift}
\leanok
\uses{def:LBTerm, def:LBTerm-mkApps, def:LBTerm-shift, def:FixDef}
$\code{NoBox}\ t$ says $t$ contains no $\Box$ anywhere, by structural recursion with the three
per-list helpers the nested list occurrences force. Box-freedom of an application spine is
box-freedom of the head and of every argument, and shifting introduces no box. The Lean
definition \code{NoBox} is by recursion on the term and lives in \code{Prop}, with the three
helper predicates and their \code{\_iff} reshapings into list quantifiers.
\end{definition}

\begin{remark}
The doc-comment cites Letouzey's box discipline [L]; this is the box-freedom a first-order
answer's conclusion asserts. It is the target-side predicate the capstone concludes of the
\emph{lowered} value; see Lemma~\ref{lem:noBox_lower_needs_noFix} for why it is not obtained by
transport along the pass.
\end{remark}

\begin{definition}[The two shapes the eraser's own output satisfies]
\label{def:NoFix}
\lean{LeanToLambdaBox.NoFix, LeanToLambdaBox.NoFixAlts, LeanToLambdaBox.NoFixAlts_iff,
LeanToLambdaBox.NoBlock, LeanToLambdaBox.NoBlockAlts, LeanToLambdaBox.NoBlockDefs,
LeanToLambdaBox.NoBlockAlts_iff, LeanToLambdaBox.NoBlockDefs_iff}
\leanok
\uses{def:LBTerm, def:FixDef}
$\code{NoFix}\ t$ says $t$ contains no \code{fix} node in spine position. A constructor node is
\emph{opaque} to it --- the clause is trivially true --- because applied-form constructor spines
carry their arguments through application and the predicate reaches them that way; \code{case}
and projection are deliberately not opaque, since a fixpoint hidden under either would satisfy
the predicate and could still take an unfolding step. $\code{NoBlock}\ t$ says no
\code{construct} node carries arguments: every constructor application is in applied (spine)
form rather than in MetaCoq's block form. Unlike the first predicate, its \code{fix} clause is
not opaque, because unfolding substitutes the block's own nodes into the body. Both Lean
definitions are structural recursions in \code{Prop} with their per-list helpers.
\end{definition}

\begin{remark}
\code{NoFix} is what the cold-start run induction carries about every term the shipping
term-level eraser returns: it never emits a fixpoint node, only the block-level pass does.
\code{NoBlock} is the applied-form invariant, the per-term predicate of the \code{ctorApplied}
clause of Definition~\ref{def:LBWfPeregrine}; measured on the emitted corpus, all $982$ emitted
constructor nodes carry an empty argument list (shipping finding F-ETA2). The Lean frontend thus
emits applied form, unlike agda2lambox and unlike MetaCoq's block convention.
\end{remark}

\begin{lemma}[The shapes survive binder closing, and the panic value]
\label{lem:noFix_toBvar}
\lean{LeanToLambdaBox.noFix_toBvar, LeanToLambdaBox.noBlock_toBvar,
LeanToLambdaBox.noFix_foldl_toBvar, LeanToLambdaBox.noBlock_foldl_toBvar,
LeanToLambdaBox.noBlock_foldl_zipIdx_map, LeanToLambdaBox.noFix_default,
LeanToLambdaBox.lbClosed_default, LeanToLambdaBox.noBlock_default}
\leanok
\uses{def:NoFix, def:toBvar, def:LBClosed}
Closing a free variable into a de Bruijn binder preserves both shapes, as do the fold forms in
which the eraser closes a body over several free variables at once --- including the
block-closing fold, which indexes its binders by name through the run's fixvar map. Moreover the
erasure's panic value is the box, which is fixpoint-free, closed at every level and invisible to
the applied-form predicate. The Lean statements for the folds are stated at the exact fold the
eraser performs, including the indexed one over a reversed zip.
\begin{proof}
\leanok
\uses{def:NoFix, def:toBvar}
Structural induction on the term. For applied form the reason is sharper than generic
congruence: abstraction maps a constructor node to a constructor node with the argument list
mapped, and that map preserves emptiness, so the single node the predicate forbids is neither
created nor destroyed. The fold forms iterate the single form; the applied-form predicate
carries no level, so the indexed fold needs no arithmetic. The panic statements are
definitional, the erasure's default value being the box.
\end{proof}
\end{lemma}

\begin{remark}
This is the honest reading of code whose ``impossible'' branches are reachable in the model:
every destructuring helper and every unreachable arm of the erasure family panics, and a panic
\emph{succeeds} in the erasure monad, returning the default. The shape induction's panic arms
are therefore discharged, not refuted.
\end{remark}

\section{The emitted program's well-formedness}
\label{sec:lowering:wf}

\begin{definition}[The subterm relation]
\label{def:SubTerm}
\lean{LeanToLambdaBox.SubTerm}
\leanok
\uses{def:LBTerm}
$\code{SubTerm}\ d\ t$ says $d$ occurs in $t$: the reflexive-transitive subterm relation. Since
\lbox{} has no binder \emph{types}, every constructor descends into every term component. The
Lean declaration \code{SubTerm} is an inductive relation with one arm per descent, the list
descents guarded by membership.
\end{definition}

\begin{definition}[Constructor and fixpoint spines, counted at maximal depth]
\label{def:ConstructSpine}
\lean{LeanToLambdaBox.IsConstructSpine, LeanToLambdaBox.CtorHeaded,
LeanToLambdaBox.ConstructSpine, LeanToLambdaBox.IsFixSpine, LeanToLambdaBox.FixHeaded,
LeanToLambdaBox.FixSpine}
\leanok
\uses{def:LBTerm, def:InductiveId, def:FixDef}
$\code{IsConstructSpine}\ t\ iid\ k\ n$ says $t$ \emph{is} a constructor spine of $(iid,k)$ with
$n$ arguments, counting both the arguments stored inside the node and those applied to it; since
this eraser emits applied form, $n$ is in practice the length of the application chain.
\code{CtorHeaded} is its existential closure. $\code{ConstructSpine}\ t\ iid\ k\ n$ is a
constructor spine \emph{occurring} in $t$, counted at its \textbf{maximal} application depth:
the application-function clause is guarded by the head not being constructor-headed, so a spine
is counted once, at its outermost application, and not again with fewer arguments --- otherwise
the saturation clause below would be unsatisfiable for any spine longer than the arity. The
three fixpoint analogues \code{IsFixSpine}, \code{FixHeaded} and \code{FixSpine} are identical
in shape at a \code{fix} head.
\end{definition}

\begin{remark}
The maximal-depth guard is the whole design point, and it is what forces the decision procedure
of Definition~\ref{def:ctorSatB} to use two mutually recursive walks.
\end{remark}

\begin{definition}[Environment and term together]
\label{def:OnProgram}
\lean{LeanToLambdaBox.OnProgram}
\leanok
\uses{def:GlobalDeclarations, def:LBTerm-envLookup}
$\code{OnProgram}\ \Gamma\ t\ P$ says the clause $P$ holds of the emitted term and of every
constant body of the emitted environment --- the environment/term split MetaRocq's
\code{expanded\_eprogram\_cstrs} makes. In Lean it is the conjunction of $P\ t$ with a bounded
quantifier over the environment's constant declarations that have a body; eight of the twelve
clauses below have this shape.
\end{definition}

\begin{definition}[The per-term clauses]
\label{def:NoDangling}
\lean{LeanToLambdaBox.NoDangling, LeanToLambdaBox.CtorsDeclared, LeanToLambdaBox.inductiveBody,
LeanToLambdaBox.CasesExhaustive, LeanToLambdaBox.ProjsDeclared, LeanToLambdaBox.FixLambda}
\leanok
\uses{def:SubTerm, def:constructorArity, def:OneInductiveBody, def:InductiveId,
def:LBTerm-envLookup}
Five per-term predicates, each quantified over the subterms of $t$. Every constant named in $t$
is declared; every constructor named in $t$ is declared, so that its arity resolves; every
\code{case} node has exactly one alternative per declared constructor of its inductive, with
alternative $i$ binding exactly the $i$-th constructor's field count; every projection selects an
in-range field of a \emph{single-constructor} inductive; and every definition of every mutual
fixpoint block has a lambda-headed body. \code{inductiveBody} is the lookup chain an inductive
identifier selects, and it is what the middle three predicates are phrased through.
\end{definition}

\begin{remark}
The exhaustiveness clause is what peregrine's \code{validate} reads and what the target's iota
rule needs; the lambda-headedness clause is coverage restriction N22, carried on the pass side as
the block bundle's \code{hfl} field.
\end{remark}

\begin{definition}[Printable binder names]
\label{def:PrintableBinders}
\lean{LeanToLambdaBox.PrintableBinderName, LeanToLambdaBox.PrintableBinders}
\leanok
\uses{def:SubTerm, def:BinderName}
A binder name is printable when it is anonymous, or is a string containing neither a double
quote nor a backslash. The printer wraps the name in a quoted atom and escapes nothing, and
peregrine's reader accepts any string atom, so the condition is exactly that the name contains
neither of the two characters that would close or escape the atom. \code{PrintableBinders} is
the conjunction of that condition at the four binder positions of a \lbox{} term --- lambda,
let, each field name of each \code{case} alternative, and each fixpoint definition's name ---
each quantified over subterms.
\end{definition}

\begin{lemma}[The alphanumeric condition is refuted on a real binder name]
\label{lem:hygienic_binder_not_alphanum}
\lean{LeanToLambdaBox.hygienic_binder_not_alphanum}
\leanok
\uses{def:PrintableBinders}
The character class the kername sanitiser establishes on a kername \emph{identifier} is false at
a binder name the eraser actually emits: a hygienic macro-scoped name contains characters
outside the alphanumeric-plus-underscore class. The Lean statement negates that class at the
literal string, and is proved by kernel evaluation of the decidable proposition.
\begin{proof}
\leanok
Kernel evaluation of the decidable proposition on the literal string.
\end{proof}
\end{lemma}

\begin{remark}
This refutation is the in-code record of a design decision. An earlier version of the
well-formedness structure carried a clause named \code{asciiNames} demanding the alphanumeric
class of every binder name. Measurement showed it false at five of the eight rungs --- one
offending binder name at three of them and thirty-four at the two largest --- and, worse, it was
not what peregrine's reader requires. The clause was restated at the condition the quoted atom
actually needs and renamed \code{printableNames}; the old field name exists nowhere in the Lean
tree.
\end{remark}

\begin{definition}[The output boundary]
\label{def:LBWfPeregrine}
\lean{LeanToLambdaBox.LBWfPeregrine}
\leanok
\uses{def:OnProgram, def:NoDangling, def:NoFix, def:ConstructSpine, def:PrintableBinders,
def:LBClosed, def:constructorArity, def:Kername}
What peregrine's untyped transform pipeline needs of the emitted program $(\Gamma, t)$, stated
on the emitted program alone, as twelve clauses: (1)~\code{keys}, no kername is declared twice;
(2)~\code{declsWf}, every inductive declaration has at least one body and every body at least
one constructor; (3)~\code{closed}, every term of the program is closed; (4)~\code{constsOk},
every constant reference resolves; (5)~\code{ctorApplied}, no constructor node carries arguments;
(6)~\code{ctorDecl}, every constructor used is declared; (7)~\code{casesExh}, every \code{case}
is exhaustive at its inductive with the right field counts; (8)~\code{fixLambda}, every fixpoint
definition is lambda-headed; (9)~\code{projDecl}, every projection is an in-range field of a
single-constructor inductive; (10)~\code{etaCtorsEnv} and (11)~\code{etaCtorsTm}, every
constructor spine, in a constant body and in the term, carries at least its constructor's
declared arity; (12)~\code{printableNames}, every binder name is printable. Clauses (3)--(9) and
(12) are stated through the environment/term combinator; the two saturation clauses are spelled
out separately because they quantify over the spine data as well.
\end{definition}

\begin{remark}
Everything but the two saturation clauses and printability is \code{peregrine validate}'s check;
the saturation invariant is what \code{remove\_params\_optimization} consumes and
\code{validate} omits. The structure mirrors MetaRocq's \code{expanded\_eprogram\_cstrs} [S,
\S7.4]. An earlier fixpoint-eta clause was false on every emitted program and was removed. At
every current rung the emitted term is a bare constant, so the \emph{term} halves of clauses
(3)--(9), (11) and (12) are vacuous and all content comes from the environment half; the
\code{etaCtorsTm} field's own docstring says so.
\end{remark}

\begin{lemma}[Reading the lambda-headedness clause at one block]
\label{lem:FixLambda-of_onProgram}
\lean{LeanToLambdaBox.FixLambda.of_onProgram}
\leanok
\uses{def:OnProgram, def:NoDangling, def:SubTerm, def:atomValue}
At any fixpoint node inside the subterm closure of a program satisfying the lambda-headedness
clause, every member of the block is lambda-headed --- stated in the indexed Boolean form the
pass relation's bundle wants. A block stored in the environment rather than reached from the
term is handled by first taking the environment half. The Lean statement takes the subterm
witness and a separate equation fixing its shape.
\begin{proof}
\leanok
\uses{def:OnProgram, def:NoDangling, def:SubTerm}
Apply the clause at the given subterm and rewrite the existential lambda form into the Boolean
head test at each index.
\end{proof}
\end{lemma}

\begin{definition}[The clause the output fails, and the true precondition]
\label{def:LBExpandedFix}
\lean{LeanToLambdaBox.LBExpandedFix, LeanToLambdaBox.PeregrinePre}
\leanok
\uses{def:OnProgram, def:ConstructSpine, def:LBWfPeregrine}
The fixpoint clause of MetaRocq's eta-expansion invariant, over environment and term: every
fixpoint occurs \emph{applied}, to a non-empty argument list longer than its principal argument
index. \code{PeregrinePre} is the conjunction of this with Definition~\ref{def:LBWfPeregrine} ---
what peregrine's first pass actually requires of its input. In Lean both are plain definitions,
the first stated through the fixpoint-spine relation of Definition~\ref{def:ConstructSpine}.
\end{definition}

\begin{remark}
This is the chapter's most important honesty statement. The clause is \emph{false} on this
eraser's output, which emits bare fixpoint constant bodies --- fifty fixpoint nodes measured,
all fifty a whole constant body, shipping finding F-ETA. Rather than weakening it or dropping it
silently, the development states it separately, names the true precondition, and has the
capstone conclude only the weaker conjunct. The gap is not paperwork: peregrine's
guarded-to-unguarded fixpoint pass is the identity on terms and discharges its whole
evaluation-preservation obligation from exactly this clause.
\end{remark}

\section{Reachability and body-less constants}
\label{sec:lowering:reach}

\begin{definition}[The delta-reachability closure]
\label{def:ReachableFrom}
\lean{LeanToLambdaBox.constRefs, LeanToLambdaBox.constRefsArgs, LeanToLambdaBox.constRefsAlts,
LeanToLambdaBox.constRefsDefs, LeanToLambdaBox.kernameElem, LeanToLambdaBox.addNames,
LeanToLambdaBox.expandStep, LeanToLambdaBox.expandRefs, LeanToLambdaBox.reachFrom,
LeanToLambdaBox.reachRefs, LeanToLambdaBox.ReachableFrom}
\leanok
\uses{def:LBTerm, def:Kername, def:GlobalDeclarations, def:LBTerm-envLookup}
$\code{constRefs}\ t$ collects every kername a term names --- the constants it references
\emph{and} the inductive block that every constructor, \code{case} and projection node reads,
since the arity and propositionality queries answer out of that declaration, so a program
reaching such a node reaches its block. One delta-step of the closure unfolds the body of every
name already seen and adds the names it references; iterating that step with fuel $|\Gamma|$
from the seed $\code{constRefs}\ t$ defines $\code{ReachableFrom}\ \Gamma\ t\ kn$. Because it is
\emph{defined} as a list-bounded closure it is decidable by construction, and the file supplies
the instance.
\end{definition}

\begin{remark}
This is the quantification domain of the environment relation's clauses, of the
specification-environment predicate's clauses, of the capstone's dependency premise and of the
body-less-reference predicate below. That $|\Gamma|$ steps really saturate is the separate
theorem that follows; without it the fixed fuel would not compose.
\end{remark}

\begin{theorem}[The fixed fuel saturates]
\label{thm:reachFrom_saturated}
\lean{LeanToLambdaBox.reachFrom_saturated, LeanToLambdaBox.find_stall,
LeanToLambdaBox.expandRefs_stall, LeanToLambdaBox.reachFrom_stall, LeanToLambdaBox.HasBody,
LeanToLambdaBox.HasBody.mem_keys, LeanToLambdaBox.envLookup_mem,
LeanToLambdaBox.mem_expandRefs}
\leanok
\uses{def:ReachableFrom}
$|\Gamma|$ delta-steps saturate the closure of any seed: beyond that nothing new is reachable.
The Lean statement \code{reachFrom\_saturated} is a list inclusion, and it rests on a one-step
characterisation \code{mem\_expandRefs} of what the expansion adds.
\begin{proof}
\leanok
\uses{def:ReachableFrom}
First, a step that unfolds no body the previous step had not already unfolded is a fixed point,
and a fixed point stays one. Then a counting argument: every step that is not already a fixed
point consumes one new key with a body, and a successful lookup exhibits the entry it answered
with, so a constant with a body is one of the environment's $|\Gamma|$ keys; a list of the
remaining such keys therefore bounds the search for a fixed point. The proof uses classical
decidability of kername equality locally, which is why the measured footprint of this branch
carries \code{Classical.choice}.
\end{proof}
\end{theorem}

\begin{lemma}[Reachability threads through the simulation]
\label{lem:ReachableFrom-through_body}
\lean{LeanToLambdaBox.ReachableFrom.subterm, LeanToLambdaBox.ReachableFrom.app,
LeanToLambdaBox.ReachableFrom.alt, LeanToLambdaBox.ReachableFrom.through_body,
LeanToLambdaBox.ReachableFrom.substList, LeanToLambdaBox.SubTerm.constRefs_subset,
LeanToLambdaBox.reachFrom_add}
\leanok
\uses{def:ReachableFrom, def:SubTerm, def:LBTerm-subst}
Three threading laws, one per shape the simulation's induction meets. What a \emph{part} of a
term reaches, the whole term reaches, a subterm naming no kername the whole term does not. What
a delta-unfolded body reaches, the program reaches --- the closure composes, because $|\Gamma|$
steps saturate it. And what an iota reduct or a beta contractum reaches, the term and the
substituends reach. The Lean statement for substitution is a disjunction, since a reachable name
may come from either side.
\begin{proof}
\leanok
\uses{def:ReachableFrom, thm:reachFrom_saturated, def:SubTerm}
The subterm law is the seed inclusion, proved by induction on the subterm relation, plus
monotonicity of the closure. The body law composes two closures into one of their sum and then
cuts back with the saturation theorem. The substitution law is the corresponding fact about
which kernames a substitution can name, lifted to the closure.
\end{proof}
\end{lemma}

\begin{remark}
These are the direction the environment-relation clauses need, since those clauses are
universally quantified over what the \emph{program} reaches and are consumed at subterms. This is
the repository's computed analogue of the side conditions MetaCoq threads through its global
erasure relation.
\end{remark}

\begin{definition}[The \code{axiom\_free} analogue]
\label{def:NoBodylessRefs}
\lean{LeanToLambdaBox.NoBodylessRefs, LeanToLambdaBox.isBodylessConst,
LeanToLambdaBox.noBodylessRefsB, LeanToLambdaBox.noBodylessRefs_iff}
\leanok
\uses{def:ReachableFrom}
No constant reachable from the emitted term is declared without a body --- [S, \S7.3]'s
\code{axiom\_free}, with no realizer whitelist. Since the closure is a list, the predicate is
decided by a list traversal, and the adequacy of that decision is proved as
\code{noBodylessRefs\_iff}, from which the file derives the instance.
\end{definition}

\begin{remark}
This predicate is load-bearing and is nothing else's premise: a run that reaches a body-less
constant is stuck at its delta step, so without this premise such a rung would be
\emph{vacuously} green. It is discharged per rung by kernel computation, and it is known to fail
on one real program --- the Fannkuch benchmark, whose emitted program declares \code{Eq.rec}
body-less (finding F-EQREC). It illustrates the repository's general method: where a premise
could hide vacuity, make it decidable and evaluate it on the actual artefacts.
\end{remark}

\section{Deciding the output boundary}
\label{sec:lowering:check}

\begin{definition}[Two generic engines]
\label{def:onProgramB}
\lean{LeanToLambdaBox.subtermAll, LeanToLambdaBox.subtermAllArgs,
LeanToLambdaBox.subtermAllAlts, LeanToLambdaBox.subtermAllDefs, LeanToLambdaBox.bodiesAllB,
LeanToLambdaBox.progB, LeanToLambdaBox.onProgramB}
\leanok
\uses{def:OnProgram, def:LBTerm, def:GlobalDeclarations}
$\code{subtermAll}\ p$ is the structural walk checking a per-node Boolean at a term and at every
subterm of it. $\code{progB}\ \Gamma\ t\ f$ checks a whole-term Boolean at the emitted term and
at every constant body --- the environment/term combinator, decided --- and \code{onProgramB} is
its composition with the walk, a per-node Boolean at every subterm of every term of the program.
All of these are Boolean-valued and structurally recursive.
\end{definition}

\begin{remark}
A constraint runs through the whole checker module and is stated in its header: every definition
must be \emph{structurally} recursive, because a well-founded definition does not reduce in the
kernel and a kernel-evaluated decision would report the instance stuck rather than failing.
\end{remark}

\begin{lemma}[From a computed check to a stated clause]
\label{lem:onProgramB_sound}
\lean{LeanToLambdaBox.subtermAll_sound, LeanToLambdaBox.progB_sound,
LeanToLambdaBox.onProgramB_sound, LeanToLambdaBox.bodiesAllB_sound}
\leanok
\uses{def:onProgramB, def:OnProgram, def:SubTerm}
The walk is exactly the subterm relation: what it checks, it checks at every subterm. Hence the
two introduction rules for the environment/term combinator --- from a whole-term check, and from
a per-node check, the latter taking as input the one place each individual clause's own
statement enters. The Lean statements are one-directional (soundness of the check, not
completeness), which is all the rungs need.
\begin{proof}
\leanok
\uses{def:onProgramB, def:SubTerm}
Induction on the subterm derivation for the walk, with the three list-descent helpers; the two
combinator rules then split the conjunction of the term check and the body check and apply the
supplied per-clause implication.
\end{proof}
\end{lemma}

\begin{definition}[Deciding saturation]
\label{def:ctorSatB}
\lean{LeanToLambdaBox.spineInfo, LeanToLambdaBox.spineOkB, LeanToLambdaBox.ctorSatB,
LeanToLambdaBox.ctorSatBFn, LeanToLambdaBox.ctorSatBArgs, LeanToLambdaBox.ctorSatBAlts,
LeanToLambdaBox.ctorSatBDefs}
\leanok
\uses{def:ConstructSpine, def:constructorArity}
\code{spineInfo} computes a node's constructor-spine data if it has any, and \code{spineOkB} is
the arity check at that node. The saturation check itself is \emph{two} mutually recursive
walks: the first checks a spine at each node, and the second walks the \textbf{function side} of
an application without treating a constructor found there as a spine root, because such a
constructor heads a longer spine whose arity is checked at the outermost application. The pair
is precisely the computational reading of the maximal-depth guard of
Definition~\ref{def:ConstructSpine}; without it the check would demand full arity at every
prefix of a spine and reject every applied-form program.
\end{definition}

\begin{lemma}[Saturation, decided]
\label{lem:ctorSat_sound}
\lean{LeanToLambdaBox.ctorSat_sound, LeanToLambdaBox.ctorSatB_root,
LeanToLambdaBox.spineOkB_sound, LeanToLambdaBox.spineInfo_of_isConstructSpine,
LeanToLambdaBox.IsConstructSpine.ctorHeaded}
\leanok
\uses{def:ctorSatB, def:ConstructSpine}
If the check succeeds on a term then every constructor spine occurring in it, at its maximal
application depth, carries at least its constructor's declared arity --- and the same for the
function-side walk. The Lean statement \code{ctorSat\_sound} is the \emph{conjunction} of the
two readings, which is what makes the induction go through.
\begin{proof}
\leanok
\uses{def:ctorSatB, def:ConstructSpine}
Induction on the spine-occurrence derivation with the conjunction of the two walks as the
motive, which is what makes the function-side clause available exactly where the guard removes
the spine-root reading. The root case reads the computed spine data off the node.
\end{proof}
\end{lemma}

\begin{remark}
This is the one non-mechanical argument in the checker file; six of the twelve clauses go
through the generic per-node engine untouched.
\end{remark}

\begin{definition}[The Boolean checker]
\label{def:lbWfPeregrineB}
\lean{LeanToLambdaBox.lbWfPeregrineB, LeanToLambdaBox.keysDistinctB, LeanToLambdaBox.declsWfB,
LeanToLambdaBox.lbClosedB, LeanToLambdaBox.noBlockB, LeanToLambdaBox.noDanglingNode,
LeanToLambdaBox.ctorsDeclaredNode, LeanToLambdaBox.altsMatchB, LeanToLambdaBox.casesExhNode,
LeanToLambdaBox.fixLambdaNode, LeanToLambdaBox.projDeclNode, LeanToLambdaBox.printableNameB,
LeanToLambdaBox.printableNode}
\leanok
\uses{def:onProgramB, def:ctorSatB, def:LBWfPeregrine}
Definition~\ref{def:LBWfPeregrine}, decided: an eleven-way conjunction for its twelve clauses,
the last saturation conjunct carrying both halves. Each conjunct is either a direct list check
(distinct keys, well-formed declarations), a whole-term check through the first engine
(closedness, applied form, saturation), or a per-node check through the second (dangling
constants, declared constructors, exhaustive cases, lambda-headed fixpoints, in-range
projections, printable names).
\end{definition}

\begin{remark}
Each node predicate's own soundness lemma is the only place that clause's statement appears; the
module header notes that plain decidability cannot reach the declaration clause, which
quantifies over an environment lookup and carries no instance --- that is why the Boolean twin
exists at all.
\end{remark}

\begin{theorem}[A kernel verdict is the well-formedness premise]
\label{thm:lbWfPeregrine_of_check}
\lean{LeanToLambdaBox.lbWfPeregrine_of_check, LeanToLambdaBox.keysDistinctB_sound,
LeanToLambdaBox.declsWfB_sound, LeanToLambdaBox.lbClosedB_sound,
LeanToLambdaBox.noBlockB_sound, LeanToLambdaBox.noDangling_of_nodes,
LeanToLambdaBox.ctorsDeclared_of_nodes, LeanToLambdaBox.casesExhaustive_of_nodes,
LeanToLambdaBox.fixLambda_of_nodes, LeanToLambdaBox.projsDeclared_of_nodes,
LeanToLambdaBox.printableBinders_of_nodes, LeanToLambdaBox.altsMatchB_sound,
LeanToLambdaBox.printableBinderName_of_check}
\leanok
\uses{def:lbWfPeregrineB, def:LBWfPeregrine}
If the checker returns true on a concrete emitted program, that program satisfies
Definition~\ref{def:LBWfPeregrine}. The Lean statement \code{lbWfPeregrine\_of\_check} takes the
Boolean equation as its only hypothesis, so a \code{decide} verdict discharges it.
\begin{proof}
\leanok
\uses{def:lbWfPeregrineB, lem:onProgramB_sound, lem:ctorSat_sound, def:LBWfPeregrine,
def:NoDangling, def:PrintableBinders, def:NoFix}
Split the eleven-way conjunction and build the structure field by field: each field is one
engine's introduction rule applied to that clause's own node soundness lemma. The saturation
conjunct is handled once, through the whole-term engine at the conjunctive statement the spine
lemma proves, and then projected to the environment half and the term half.
\end{proof}
\end{theorem}

\begin{remark}
The measured axiom footprint is \code{propext}, \code{Classical.choice} and \code{Quot.sound}
--- no \code{sorryAx} --- and the row is a committed fixture in \code{test/ledger.expected}.
This is how \emph{every} green rung supplies the capstone's well-formedness premise, each rung
literally applying this theorem to a kernel-evaluated decision; the premise is discharged, not
assumed. Two anonymous examples close the file --- one exhibits a concrete accepted
two-declaration program, the other a rejected one, namely the same term at an empty environment
--- so the checker is neither vacuous nor constantly true. Having no Lean names, they are not
nodes of this graph.
\end{remark}

\section{The verified pass template: \texorpdfstring{\code{optimize}}{optimize}}
\label{sec:lowering:optimize}

\begin{definition}[The propositional-case optimisation]
\label{def:LBOptimize}
\lean{LeanToLambdaBox.caseCollapse, LeanToLambdaBox.projCollapse, LeanToLambdaBox.LBOptimize,
LeanToLambdaBox.LBOptimizeArgs, LeanToLambdaBox.LBOptimizeAlts, LeanToLambdaBox.LBOptimizeDefs,
LeanToLambdaBox.LBOptimize_env, LeanToLambdaBox.optDecl}
\leanok
\uses{def:LBTerm, def:LBTerm-subst, def:isPropositionalInductive, def:GlobalDeclarations}
MetaRocq's \code{optimize} / \code{remove\_match\_on\_box} [S, \S7.4], as a \textbf{function}. A
\code{case} on a propositional single-constructor inductive collapses to its one alternative's
body with one copy of $\Box$ per field binder simultaneously substituted; a projection out of a
propositional inductive collapses to $\Box$; everything else is structural. On environments,
every constant body is optimised with respect to the \emph{original} environment, and inductive
and body-less declarations are left untouched. The three auxiliary list functions exist because
the structural-recursion checker does not see through a list map at the nested list occurrences.
\end{definition}

\begin{remark}
This is the syntactic residue of Letouzey's rule that computationally irrelevant matches
disappear [L]. That the environment pass reads the pre-pass environment is harmless, because
propositionality is invariant under it.
\end{remark}

\begin{lemma}[The commutation package]
\label{lem:LBOptimize_subst_comm}
\lean{LeanToLambdaBox.box_subst_swap_gen, LeanToLambdaBox.substList_replicate_box_subst,
LeanToLambdaBox.box_subst_shift_swap, LeanToLambdaBox.substList_replicate_box_shift,
LeanToLambdaBox.caseCollapse_subst, LeanToLambdaBox.caseCollapse_shift,
LeanToLambdaBox.LBOptimize_shift_comm, LeanToLambdaBox.LBOptimize_subst_comm,
LeanToLambdaBox.LBOptimize_substList_box, LeanToLambdaBox.LBOptimize_substList,
LeanToLambdaBox.LBOptimize_mkApps, LeanToLambdaBox.LBOptimize_iota_red,
LeanToLambdaBox.LBOptimize_fixSubst, LeanToLambdaBox.LBOptimize_fixUnfold_body,
LeanToLambdaBox.LBOptimize_case_noncollapse, LeanToLambdaBox.LBOptimize_case_collapse,
LeanToLambdaBox.LBOptimize_proj_noncollapse, LeanToLambdaBox.LBOptimize_proj_collapse,
LeanToLambdaBox.envLookup_LBOptimize_env,
LeanToLambdaBox.isPropositionalInductive_LBOptimize_env,
LeanToLambdaBox.isStuckApp_LBOptimize, LeanToLambdaBox.isStuckApp_LBOptimize_of_value}
\leanok
\uses{def:LBOptimize, def:LBTerm-subst, def:LBTerm-shift, def:LBTerm-fixSubst, def:Value,
def:isStuckApp, def:isPropositionalInductive}
Everything the correctness proof needs of the pass, all resting on one observation: $\Box$ is
closed, so substituting boxes commutes with everything. Substituting $\Box$ at depth $d$
commutes past an outer substitution, lowering its depth; iterated, it commutes past a
substitution and past a shift, lowering the depth or cutoff by the number of boxes. Hence the
collapse decision commutes with substitution and shifting, hence the pass commutes with
shifting, with substitution, with a list of boxes, with a simultaneous substitution, with an
application spine, with an iota reduction's drop-and-reverse, and with fixpoint unfolding. The
two readings of the pass at a \code{case} and at a projection are split by decidable side
conditions; looking a constant up in the optimised environment yields the optimised body;
optimising the environment does not change which inductives are propositional; and on a
\emph{value}, the pass preserves being a stuck application head.
\begin{proof}
\leanok
\uses{def:LBOptimize}
The arithmetic core is a single substitution-against-substitution swap that degenerates cleanly
because the substituted term is closed and needs no shift bookkeeping. The dangerous case
throughout is the propositional single-branch \code{case}, where the pass \emph{removes}
binders: the two readings line up exactly because the boxes substituted in are closed. The
stuck-head statement is an induction on the value relation --- a value's spine head can only be
a free variable when it is stuck, never a \code{case} or projection the pass might collapse, so
the head shape survives; the stuckness test reads only one flag, so the flag on the value need
not match the query.
\end{proof}
\end{lemma}

\begin{theorem}[The pass discharges the propositional-case flag]
\label{thm:LBOptimize_correct}
\lean{LeanToLambdaBox.LBOptimize_correct}
\leanok
\uses{def:LBOptimize, def:WcbvEval, def:Eval}
If a \lbox{} term evaluates to a value under weak call-by-value evaluation \emph{with the
propositional-case rules enabled}, then the optimised term evaluates, in the optimised
environment, to the optimised value under the same evaluation \emph{with those rules disabled}.
This is exactly MetaRocq's \code{optimize\_correct} shape: the pass discharges a flag. The Lean
statement relates the two flag points \code{propBlockFlags} and \code{blockFlags}, both of which
are \emph{block}-form constructor regimes.
\begin{proof}
\leanok
\uses{def:LBOptimize, lem:LBOptimize_subst_comm, def:WcbvEval, thm:eval_to_value}
Induction on the evaluation derivation. The two flag-gated rules --- the single-branch
propositional \code{case} and the propositional projection --- are discharged by the
corresponding collapse equation; the guarded non-propositional iota and projection rules re-use
the same rule on the optimised node; fixpoint unfolding commutes by the unfolding equation; and
the applied-form rules together with the unguarded fixpoint rule are unreachable at the source
flags and are closed at the flag hypothesis.
\end{proof}
\end{theorem}

\begin{remark}
The measured axiom footprint is \code{propext} and \code{Quot.sound}: no
\code{Classical.choice}, no \code{sorryAx}, the cleanest result in the repository. Two caveats
belong beside it, both visible in the code. First, the theorem is stated at block-form flags on
both sides, while the shipping erasure emits applied form and the capstone concludes at the
erasure flags; the proved pass correctness therefore cannot be composed with the capstone
without re-proving its arms. Second, the erasure never marks an inductive propositional, so on
emitted environments the flag is inert and the pass would be the identity. The file is outside
the import closure of the capstone and of the rungs: it is a faithful re-proof of [S, \S7.4] and
the declared template the pass relation is stated in the shape of, not a component of the
shipping chain.
\end{remark}

\begin{lemma}[Three vacuity guards]
\label{lem:LBOptimize_correct_hyps_satisfiable}
\lean{LeanToLambdaBox.LBOptimize_correct_not_refutable,
LeanToLambdaBox.LBOptimize_correct_hyps_satisfiable, LeanToLambdaBox.LBOptimize_correct_fires,
LeanToLambdaBox.vac_optimize_collapses, LeanToLambdaBox.vacTerm}
\leanok
\uses{thm:LBOptimize_correct, def:LBOptimize}
On an environment declaring one propositional single-constructor inductive with one field, and
the single-branch \code{case} on $\Box$: the hypothesis class of
Theorem~\ref{thm:LBOptimize_correct} is non-empty; this concrete term does satisfy the
hypothesis, by the flag-gated single-branch rule; the pass really does collapse it to $\Box$;
and applying the theorem at it yields genuine evaluation content rather than a trivial
statement. The Lean statements are at the fixture's explicit environment and term.
\begin{proof}
\leanok
\uses{def:LBOptimize, def:WcbvEval}
The environment and term are explicit, so the collapse is a definitional computation; the source
evaluation is one application of the flag-gated single-branch rule at the propositional
inductive; and the third statement is the theorem applied to the second.
\end{proof}
\end{lemma}

\begin{remark}
This three-part pattern --- hypothesis class inhabited, concrete witness, firing yields content
--- is the repository's standard anti-vacuity discipline, and it recurs in later chapters.
\end{remark}

\section*{Deviations from the paper and caveats}

Five points, each already carried by a node above, deserve restating together. First, the pass
is \textbf{one-directional}: the eliminator arm drops the arguments before the discriminant with
no premise about them, which is sound for a forward simulation and is what MetaRocq's own
expansion does, but the relation is not a bisimulation and must never be described as
semantics-preserving in both directions. Second, the pass is \textbf{non-deterministic} at a
block member, and Definition~\ref{def:LowerFix} hides its lowered bodies and fix variables
behind existentials, so every ``the pass emits'' in the prose is ``there is an emission such
that''. Third, the \emph{source} of this pass is a \lbox{} dialect in which eliminators are
still constants and recursion is still by environment reference, so the pass manufactures the
two node kinds MetaRocq's erasure gets for free from Rocq's primitive \code{match} and
\code{fix} --- which is exactly why the resulting fixpoint node is not inert under box transport
(Lemma~\ref{lem:noBox_lower_needs_noFix}). Fourth, the emitted output is in \textbf{applied}
constructor form, not MetaCoq's block form, and the shipped evaluation flags say so, whereas the
one fully verified pass in the tree, Theorem~\ref{thm:LBOptimize_correct}, is stated at
block-form flags. Fifth, the conclusion about the emitted program is
Definition~\ref{def:LBWfPeregrine} and explicitly \emph{not} \code{PeregrinePre}: the fixpoint
eta clause is false on this eraser's output, and that gap is a named shipping finding rather
than paperwork.
